\section{Introducción}

La evolución de los sistemas electrónicos en el sector de la automoción ha permitido que los vehículos modernos cuenten con redes de comunicación internas avanzadas, como el protocolo \textit{CAN BUS}, que interconectan múltiples unidades de control electrónico (ECUs) para funciones de confort, seguridad, motorización y multimedia. Sin embargo, muchos vehículos fabricados en décadas anteriores carecen de estas capacidades, lo que limita significativamente su nivel de automatización y conectividad.\\

Este proyecto surge con el objetivo de transformar un vehículo antiguo —equipado originalmente únicamente con una unidad de control de motor conectada a través del conector OBD para diagnóstico— en una plataforma más avanzada, dotada de nuevas funcionalidades que permitan ampliar su experiencia de uso y control interno.\\

Para ello, se propone la integración de un sistema de automatización distribuido basado en el protocolo \textit{CAN BUS}, empleando controladores diseñados por el propio equipo de desarrollo. Estos controladores estarán basados en el chip \textbf{MCP2515}, un controlador CAN SPI ampliamente utilizado, que permitirá establecer nodos adicionales dentro de la red del vehículo.\\

El proyecto contempla la creación de al menos dos nuevas unidades funcionales:

\begin{itemize}
  \item \textbf{Unidad de confort:} encargada de funciones como el control de climatización, cierre centralizado, control de ventanas y posibles ajustes de iluminación interior.
  \item \textbf{Unidad multimedia:} enfocada en la gestión de audio, interfaces con el usuario, conectividad y notificaciones del sistema.
\end{itemize}

Estas unidades estarán interconectadas mediante un nuevo bus CAN estructurado, que permitirá comunicación en tiempo real y expansión futura. El sistema será desarrollado, probado e integrado completamente por el equipo, aprovechando tanto hardware propio como soluciones abiertas adaptadas al entorno del vehículo.\\

Con esta iniciativa se busca no solo modernizar un automóvil existente, sino también explorar las posibilidades de extender tecnologías de automoción a plataformas donde originalmente no estaban concebidas, lo cual abre nuevas puertas para la actualización tecnológica de flotas antiguas o proyectos de restauración inteligente.

